\section{Constante d'Euler--Kronecker du caractère dodécaphasique}

La factorisation de Dedekind

\begin{equation}
\zeta_{\mathbb Q(\sqrt3)}(s)=\zeta(s)L(s,\chi_{12})
\label{eq:dedekind-sqrt3}
\end{equation}

donne, en comparant les termes constants en \(s=1\),

\begin{equation}
\gamma_{\mathbb Q(\sqrt3)}
=\gamma+\frac{L'(1,\chi_{12})}{L(1,\chi_{12})}
=1.0539656082484825800\ldots.
\label{eq:euler-kronecker}
\end{equation}

Cette constante représente le terme fini normalisé au pôle de la
fonction zêta du corps réel dodécaphasique. Elle ne constitue pas une variante de
\(\gamma\) d'Euler, mais la correction arithmétique induite par
\(\chi_{12}\).

En effet, en écrivant \(u=s-1\),
\[
\zeta(s)=u^{-1}+\gamma+O(u),
\qquad
L(s,\chi_{12})=L_1+L_1'u+O(u^2),
\]
on obtient
\begin{equation}
\zeta_{\mathbb Q(\sqrt3)}(s)
=\frac{L_1}{u}+\bigl(\gamma L_1+L_1'\bigr)+O(u).
\label{eq:ek-laurent-derivation}
\end{equation}
La constante d'Euler--Kronecker est le quotient du terme fini par
le résidu ; diviser \cref{eq:ek-laurent-derivation} par \(L_1\) prouve
\cref{eq:euler-kronecker}. Cette dérivation explique simultanément le
résidu
\[
\operatorname*{Res}_{s=1}\zeta_{\mathbb Q(\sqrt3)}(s)
=\frac{\log(2+\sqrt3)}{\sqrt3}
\]
et la correction finie. La valeur numérique est certifiée en évaluant
\(L(1,\chi_{12})\) et \(L'(1,\chi_{12})\) avec le même tableau résiduel et une
borne de queue ; elle n'est pas importée d'un tableau de corps quadratiques.

\section{Constantes arithmétiques associées aux nombres premiers}

Une fois les cycles arithmétiques irréductibles sélectionnés, on obtient deux
familles distinctes :

\begin{align}
B_1&=\lim_{x\to\infty}\left(
\sum_{p\le x}\frac1p-\log\log x
\right),
\label{eq:meissel-mertens}\\
C_2&=\prod_{p>2}\frac{p(p-2)}{(p-1)^2}.
\label{eq:twin-prime-constant}
\end{align}

\(B_1\) régularise une somme locale ; \(C_2\) est une densité multiplicative de paires jumelles. Les deux exigent un contrôle de queue. Ils ne se déduisent pas l'un de l'autre parce qu'ils partagent les nombres premiers.

